\chapter{More on Loops and Iteration in Python}

\section{Introduction: Real-Life Hook}

Students face repeated tasks in daily life:  

\begin{itemize}
    \item Checking homework for all subjects  
    \item Sending the same message to multiple friends  
    \item Calculating total marks for all students in a class  
\end{itemize}

In Python, loops help us \textbf{automate repeated tasks efficiently}, especially when working with \textbf{lists, strings, and dictionaries}.

\section{Concept: Looping Through Lists, Strings, and Dictionaries}

\subsection{Looping Through Lists}

\begin{verbatim}
fruits = ["apple", "banana", "mango"]
for fruit in fruits:
    print("I like", fruit)
\end{verbatim}

- Loops repeat an action for each element in the list  

\subsection{Looping Through Strings}

\begin{verbatim}
word = "Python"
for char in word:
    print(char)
\end{verbatim}

- Each character in the string is processed individually  

\subsection{Looping Through Dictionaries}

\begin{verbatim}
student = {"name": "Amina", "marks": 90, "age": 15}

# Loop through keys
for key in student:
    print(key, ":", student[key])

# Loop through key-value pairs
for key, value in student.items():
    print(key, "->", value)
\end{verbatim}

- Loops can iterate through keys or key-value pairs for efficient data handling  

\section{Nested Loops}

\begin{verbatim}
marks = [
    [80, 70, 90],
    [75, 85, 65],
    [90, 95, 100]
]

for i, student in enumerate(marks):
    print("Student", i+1, "marks:")
    for mark in student:
        print(mark, end=" ")
    print()
\end{verbatim}

- Outer loop iterates through students, inner loop through marks  
- Nested loops are powerful for \textbf{2D data}  

\section{Using \texttt{range()} with Loops}

\begin{verbatim}
# Print numbers 1 to 5
for i in range(1, 6):
    print(i)

# Print even numbers 2 to 10
for i in range(2, 11, 2):
    print(i)
\end{verbatim}

- \texttt{range(start, stop, step)} controls the loop index  

\section{Loop Control Statements}

\begin{itemize}
    \item \textbf{break} — exit loop early  
\begin{verbatim}
for i in range(5):
    if i == 3:
        break
    print(i)  # prints 0 1 2
\end{verbatim}

    \item \textbf{continue} — skip current iteration  
\begin{verbatim}
for i in range(5):
    if i == 2:
        continue
    print(i)  # prints 0 1 3 4
\end{verbatim}
\end{itemize}

\section{Mini-Project: Grade Calculator}

Create a program to calculate total and average marks for 3 students:

\begin{verbatim}
students = [
    {"name": "Amina", "marks": [80, 70, 90]},
    {"name": "Rafi", "marks": [75, 85, 65]},
    {"name": "Sadia", "marks": [90, 95, 100]}
]

for student in students:
    total = sum(student["marks"])
    avg = total / len(student["marks"])
    print(student["name"], "- Total:", total, "Average:", avg)
\end{verbatim}

Output:

\begin{verbatim}
Amina - Total: 240 Average: 80.0
Rafi - Total: 225 Average: 75.0
Sadia - Total: 285 Average: 95.0
\end{verbatim}

\section{Step-by-Step Diagram: Nested Loops with Dictionaries}

\begin{center}
\begin{tikzpicture}[node distance=1.5cm]
\node (start) [draw, ellipse, fill=green!20] {Start Program};
\node (outer) [draw, rectangle, fill=yellow!20, below of=start] {Loop Through Outer Collection (Students)};
\node (inner) [draw, rectangle, fill=orange!20, below of=outer] {Loop Through Inner Collection (Marks)};
\node (process) [draw, rectangle, fill=orange!20, below of=inner] {Process Each Element (Sum / Average)};
\node (output) [draw, ellipse, fill=green!20, below of=process] {Display Results};
\draw[->] (start) -- (outer);
\draw[->] (outer) -- (inner);
\draw[->] (inner) -- (process);
\draw[->] (process) -- (output);
\end{tikzpicture}
\end{center}

\section{Key Concepts}

\begin{itemize}
    \item Loops iterate through lists, strings, and dictionaries  
    \item Nested loops handle multi-dimensional or hierarchical data  
    \item \texttt{range()} allows index-based loops with steps  
    \item \texttt{break} and \texttt{continue} control loop flow  
\end{itemize}

\section{End-of-Chapter Exercises}

\begin{exercise}{Very Short Questions}
\begin{enumerate}
    \item What is the purpose of a loop?  
    \item How do you skip an iteration in a loop?  
    \item What is a nested loop?  
\end{enumerate}
\end{exercise}

\begin{exercise}{Short Answer Questions}
\begin{enumerate}
    \item Write a loop to print numbers 10 to 1 in reverse.  
    \item Explain the difference between \texttt{break} and \texttt{continue}.  
    \item Write a nested loop to print a 3x3 multiplication table.  
\end{enumerate}
\end{exercise}

\begin{exercise}{Activity-Based Questions}
\begin{enumerate}
    \item Ask the user to input 3 students’ names and 3 marks each. Calculate total and average using nested loops.  
    \item Create a program to print all even numbers between 1 and 20 using \texttt{continue}.  
    \item Write a program to display each character of a string in reverse order using a loop.  
\end{enumerate}
\end{exercise}

\section{Teacher Notes}

\marginnote{
\textbf{Learning Objectives:}  
\begin{itemize}
    \item Students can use loops with lists, strings, and dictionaries  
    \item Students can implement nested loops for complex data  
    \item Students can control loop flow using \texttt{break} and \texttt{continue}  
\end{itemize}

\textbf{Teaching Tips:}  
\begin{itemize}
    \item Demonstrate with real-life data like students’ marks, names, and messages  
    \item Start with simple loops, then introduce nested loops  
    \item Show \texttt{break} and \texttt{continue} in context for clarity  
\end{itemize}

\textbf{Assessment:}  
\begin{itemize}
    \item Check correct implementation of nested loops  
    \item Observe correct usage of loop control statements  
\end{itemize}
}